\documentclass[10pt,a4paper]{amsart}
\usepackage[margin=19mm]{geometry}
\usepackage{amsmath,amssymb,mathtools}
\usepackage[scr=rsfs,cal=euler]{mathalfa}
\usepackage{xfrac}
\usepackage[unicode,hidelinks,bookmarks=false]{hyperref}
\newcommand{\ie}{i.e.\ }
\newcommand{\cf}{cf.\ }
\newcommand{\resp}{respectively\ }
\newcommand{\Subsection}{\subsection}
\providecommand{\nobreakdash}{\nobreak\hbox{-}}
\setlength{\emergencystretch}{2em}
\multlinegap0pt
\usepackage{amssymb}
\usepackage{mathtools}
\newtheorem{corollary}{Corollary}
\newtheorem{lemma}{Lemma}
\newtheorem{cla}{Claim}
\newtheorem{poc}{Poc}
\title{\bf\Large Strong non-principality of positive codegree Tur\'an density}
\author{Levente Bodn\'ar, Jun Gao, Oleg Pikhurko, Mingyuan Rong, Shumin Sun}
\date{Source: arXiv:2606.20494v1 (CC BY 4.0)}
\begin{document}
\maketitle

\noindent\emph{Source and licence.} \bf\Large Strong non-principality of positive codegree Tur\'an density. Version arXiv:2606.20494v1 (CC BY 4.0), distributed by arXiv under the Creative Commons Attribution 4.0 International licence, http://creativecommons.org/licenses/by/4.0/, as recorded in the arXiv licence metadata for this exact version and shown on the abstract page https://arxiv.org/abs/2606.20494v1. Full licence notice: \texttt{licenses/arxiv-2606.20494-CC-BY-4.0.txt}.\par\smallskip

\noindent\textit{Editorial scope.} Counted unit: the article's lemma lem:main (source0.tex lines 333--335). The article's complete original proof of it is delivered with it (source0.tex lines 336--382, one \texttt{proof} environment of 47 lines). No mathematics is written, repaired or re-derived: the statement and the proof are the article's own lines, in the article's own notation, and no retained line is substituted or re-flowed.  Retained uncounted as the source-local context the proof needs: lines 283--285.  Nothing was omitted from any retained span, so the package carries no omission record.  No retained line contains a citation, so the wrapper carries no bibliography and no bibliography entry.  No retained line contains a figure environment, an image-inclusion command or drawing code, so no image is delivered and the package declares no asset dependency.  Editorial formatting only, in addition: an AMS \texttt{amsart} preamble that restates the article's own declarations and macros used by the retained lines, namely the environments \texttt{corollary}, \texttt{lemma}, \texttt{cla}, \texttt{poc} on separate counters, together with the standard packages those lines need.  The retained environments print in this document as Corollary 1, Lemma 1, Cla 1, Poc 1 in order of appearance, whereas the article prints the same items with the numbering of its own full document.  The complete native source of the article as distributed in the arXiv e-print for this exact version is bundled unchanged as \texttt{sources/203222/source0.tex}.
\begin{corollary}\label{thm:H graph}
    For any integers $r\ge 1$ and $k\ge 2$, we have that $\gamma^+(H_r^k)=\frac{r-1}{r}$ and $H_r^k$ is not $r$-colourable.
\end{corollary}

\begin{lemma}\label{lem:main}
    For any integer $k\ge 3$, we have $\frac{1}{k}\le\gamma^+(Q_2^k,D_{k+1}^k) \le \frac{1}{2}-\alpha$, where $\alpha := \frac{1}{10(k-1)}$. 
\end{lemma}

\begin{proof}
    By considering the balanced complete $k$-partite $k$-graph, it is easy to see that $ \gamma^+(Q_2^k,D_{k+1}^k) \ge \frac{1}{k}$.
    Next we will show $\gamma^+(Q_2^k,D_{k+1}^k) \le \frac{1}{2}-\alpha$.
    Let $G$ be a $k$-graph on $n$ vertices with $\delta_{k-1}^+(G)\ge\left(\frac{1}{2}-\alpha \right)n$, where $n$ is sufficiently large.
    Suppose that $G$ is $Q_2^k$-free and $D_{k+1}^k$-free.
    First, we show that $\partial_{k-1}(G)$ has large minimum positive codegree.
    %, namely larger than $(\frac{1}{2}+3\alpha)n$, by the $D_{k+1}^k$-freeness. 
    
    
    \begin{cla}\label{cl:k-2degree}
    $\delta_{k-2}^+(\partial_{k-1}(G))\ge\left(\frac{1}{2}+3\alpha\right)n$.
    \end{cla}
    \begin{poc}
    For any $(k-2)$-set $S\subseteq V(G)$, let $R$ be the link graph of $S$ in $G$, i.e., 
    \[
    \mbox{$V(R)=\{v:\{v\}\cup S\in\partial_{k-1}(G)\}$ \ \ and \ \ $E(R)=\{uv:\{u,v\}\cup S\in G\}$.}
    \] Our aim is to show that $|V(R)|\ge\left(\frac{1}{2}+3\alpha\right)n$ whenever $E(R)$ is non-empty. Suppose that $E(R)$ is non-empty. Then the minimum degree of $R$ is at least $\left(\frac{1}{2}-\alpha\right)n$. Since $G$ is $D_{k+1}^k$-free, we have that $R$ is $K_{k+1}$-free. It follows from Tur\'an's theorem that $|V(R)|\ge \frac{k}{k-1}\delta_1(R)-1\ge \left(\frac{1}{2}+3\alpha\right)n$.
    \end{poc}
    
    \begin{cla}\label{cl:IndepY}
    There exists a vertex set $Y\subseteq V(G)$ of size at least $\left(\frac{1}{2}-2\alpha\right)n$ such that each edge of $G$ intersects $Y$ in at most one vertex.
    \end{cla}
    \begin{poc}
    Suppose, for a contradiction, that every set of size at least $\left(\frac{1}{2}-2\alpha\right)n$ contains at least two vertices of some edge.  By Claim~\ref{cl:k-2degree}, the $(k-1)$-graph $\partial_{k-1}(G)$ has positive codegree at least $\left(\frac{1}{2}+3\alpha\right)n$. By Corollary~\ref{thm:H graph}, it contains $H_2^{k-1}$ as a subgraph, which we denote by $X$. Let $S_1,S_2,\dots, S_m$ be all $(k-1)$-edges of $X$.
    
    We now prove that there exist disjoint edges $F_1,\dots,F_m\subseteq V(G)\setminus V(X)$ with $F_i=\{x_1^i,\dots,x_k^i\}$ such that $S_i\cup\{x_1^i\}$ and $S_i\cup\{x_2^i\}$ are edges of $G$ for every $1\le i\le m$.
    
    Suppose that we have found $F_1,\dots,F_{i-1}$ for some $i\in[m]$. Let $A:=X\cup\bigcup_{j=1}^{i-1}F_j$. Since $\delta_{k-1}^+(G)\ge\left(\frac{1}{2}-\alpha\right)n$ and $n$ is sufficiently large, we have
    \[
    \left|N(S_i)\setminus A\right|\ge\left(\frac{1}{2}-2\alpha\right)n.
    \]
    By our assumption for contradiction, there exists an edge $F_i$ intersecting $N(S_i)\setminus A$ in at least two vertices, call them $x_1^i$ and $x_2^i$. Among such edges pick one with $|F_i\setminus A|$ as large as possible. Since $\delta_{k-1}^+(G)\ge\left(\frac{1}{2}-\alpha\right)n>|A|$, a familiar argument shows that $F_i$ must be disjoint from $A$, as desired.
    Then $F_1,\dots,F_m$ together with $X$ form a copy of $Q_2^k$, a contradiction.
   \end{poc}
   Let $Y$ be a vertex set of size at least $\left(\frac{1}{2}-2\alpha\right)n$ such that each edge of $G$ intersects $Y$ in at most one vertex.
   If no edge intersects $Y$, then for any edge $E\in E(G)$ and any $(k-2)$-subset $S\subseteq E$, we have $N_{\partial_{k-1}(G)}(S)\subseteq V(G)\setminus Y$ and
   \[
   |N_{\partial_{k-1}(G)}(S)|\le n-|Y|\le\left(\frac{1}{2}+2\alpha\right)n,
   \]
   contradicting Claim~\ref{cl:k-2degree}.
   If some edge intersects $Y$ in exactly one vertex, take such an edge $E$ and choose a $(k-2)$-subset $S\subseteq E$ with $S\cap Y\neq\emptyset$. Again $N_{\partial_{k-1}(G)}(S)\subseteq V(G)\setminus Y$, so
   \[
   |N_{\partial_{k-1}(G)}(S)|\le n-|Y|\le\left(\frac{1}{2}+2\alpha\right)n,
   \]
   contradicting Claim~\ref{cl:k-2degree}. 
   Thus $G$ contains $Q_2^k$ or $D_{k+1}^k$ as a subgraph, which finishes the proof of Lemma~\ref{lem:main}.
\end{proof}

\end{document}
